\setchapterpreamble[u]{\margintoc}
\chapter{What a Learning Problem Is}
\labch{learning-problem}

\sources{Zhang, Lipton, Li and Smola, \emph{Dive into Deep Learning}
\cite{d2l2023}, Chapter~1.}

Before any architecture is chosen, a learning problem has to be stated. This
chapter fixes the vocabulary and the four moving parts that every supervised
method rearranges, including the ones this project is preparing for.

\section{Vocabulary}
\labsec{vocabulary}

A few words are used with technical precision throughout Part~IV.

\begin{description}
	\item[Machine learning] fitting a program's behaviour from data instead of
	specifying it by hand --- \emph{programming with data}.
	\item[Deep learning] the subset in which the fitted function is a
	composition of many parameterised layers.
	\item[Dataset] the collection of recorded observations.
	\item[Model] the parameterised family of functions being fitted.
	\item[Parameter] a number adjusted during training.
	\item[Learning algorithm] the procedure that adjusts the parameters.
	\item[Training] the act of running that procedure.
\end{description}

\marginnote{Keeping ``model'' for the family and ``parameters'' for the numbers
inside it avoids a persistent ambiguity: a trained model is a family
\emph{plus} a particular setting, and the two are chosen by different means.}

\section{The four components}
\labsec{four-components}

Any learning problem, however it is later dressed up, consists of four things
that must be specified separately:

\begin{description}
	\item[Data] what the method is allowed to learn from;
	\item[Model] the family of transformations available to it;
	\item[Objective function] a number that says how badly the current model is
	doing;
	\item[Algorithm] the procedure that adjusts the model to reduce that number.
\end{description}

\marginnote{Keeping the four apart is worth the pedantry. Most confusing results
come from changing one of them silently --- a new loss, a new sampling scheme
--- while attributing the difference to the model.}

\subsection{Data}
\labsubsec{data}

A dataset is a collection of \emph{examples}, also called samples or data
points. Each example is described by a fixed list of \emph{features}, also
called covariates, and in supervised problems it carries a \emph{label} --- the
special attribute to be predicted, also called the target. The number of
features is the \emph{dimensionality} of the example, and when that number is the
same for every example the data are of fixed length, which is what makes it
possible to write a batch of \(n\) examples as a matrix \(X\in\R^{n\times d}\) and
reuse everything from \refch{learning-from-data}.

\begin{definition}
Examples are \emph{i.i.d.} --- independent and identically distributed ---
when each is drawn from the same distribution, independently of the others. Nearly every
guarantee about generalisation assumes this, and nearly every deployed system
violates it to some degree.
\end{definition}

The formal content of that definition is \refsec{independence}, and the reason it
matters is \refsec{lln}: without it, an average over the sample stops estimating
anything in particular.

\subsection{Model}
\labsubsec{model}

The model is a parameterised map from features to predictions. Fixing the model
family fixes what the method can express; fixing the parameters within that
family fixes what it currently predicts. For a linear model the reachable
predictions are exactly the column space of the data matrix, which is the
representational question of \refch{vectors-and-column-space} arriving in its
applied form.

\subsection{Objective function}
\labsubsec{objective}

The objective --- the loss or cost function --- turns the gap between prediction
and label into a single number, lower being better. For a real-valued target the
usual choice is squared error,
\[
	\ell(\hat y,y)=\tfrac12(\hat y-y)^{2},
\]
and for a categorical target it is cross-entropy. \refsec{losses-from-likelihood}
shows that these two are not conventions but the maximum likelihood objectives
for Gaussian and Bernoulli noise respectively.

\begin{remark}
The quantity one actually cares about is often not differentiable ---
classification accuracy, ranking quality, revenue. What gets optimised is a
tractable surrogate, and the gap between the surrogate and the real goal is a
permanent source of surprise rather than a temporary inconvenience.
\end{remark}

\subsection{Optimisation algorithm}
\labsubsec{algorithm}

The algorithm searches the parameter space for a setting that makes the
objective small. Gradient descent and its variants dominate, for the practical
reason established in \refsec{backprop}: the gradient of a composed
differentiable model can be computed at roughly the cost of evaluating it, no
matter how many parameters it has.

\section{Training, testing, and overfitting}
\labsec{training-testing}

The available data are split. The \emph{training set} is used to fit the
parameters; the \emph{test set} is held out and used only to estimate how the
fitted model behaves on examples it has never seen.

\marginnote{A test set inspected repeatedly stops being a test set. Each look
leaks information about it into the modelling decisions.}

The split exists because loss on the training set is a biased estimate of
quality. A model flexible enough can drive training loss towards zero by
memorising the particular examples it was given, including their noise.
\emph{Overfitting} is the name for the resulting gap: training loss keeps
falling while test loss turns and rises.

\begin{remark}
Overfitting is not a defect of a specific model family. It is what happens
whenever the capacity of the family is large relative to the information in the
data, and it is the reason that model selection, regularisation, and honest
held-out evaluation are part of the method rather than optional extras.
\end{remark}

By \refsec{clt}, a test set of \(n\) examples estimates the true error with a
standard deviation of order \(1/\sqrt n\). Reported differences smaller than that
are not evidence of anything, however carefully the experiment was run.

\section{Why the framing is worth keeping}
\labsec{framing}

The four components are independent choices, and almost every method discussed
in \refch{map-of-learning-problems} is a different setting of them rather than a
different subject. Recognising that keeps the number of things to learn small:
new architectures change the model, new losses change the objective, and new
optimisers change the algorithm, but the shape of the problem stays put.
